\documentclass[11pt,a4paper]{article}
\usepackage[utf8]{inputenc}
\usepackage[margin=2.2cm]{geometry}
\usepackage{amsmath,amssymb,amsfonts,bm}
\usepackage{booktabs}
\usepackage{graphicx}
\usepackage{xcolor}
\usepackage{hyperref}
\usepackage{listings}
\usepackage{tcolorbox}
\usepackage{algorithm}
\usepackage{algpseudocode}
\usepackage{microtype}
\usepackage{array}

\definecolor{primary}{RGB}{31, 78, 121}
\definecolor{secondary}{RGB}{0, 128, 128}
\definecolor{accent}{RGB}{180, 50, 50}
\definecolor{darkgray}{RGB}{50, 50, 50}
\definecolor{codebg}{RGB}{248, 249, 250}
\definecolor{bordercolor}{RGB}{220, 224, 230}

\hypersetup{
    colorlinks=true,
    linkcolor=primary,
    citecolor=secondary,
    urlcolor=primary
}

\lstset{
    backgroundcolor=\color{codebg},
    basicstyle=\ttfamily\small,
    breaklines=true,
    frame=single,
    rulecolor=\color{bordercolor},
    keywordstyle=\color{primary}\bfseries,
    commentstyle=\color{gray}\itshape,
    stringstyle=\color{secondary}
}

\tcbuselibrary{skins,breakable}
\newtcolorbox{infobox}[1][]{
    enhanced,
    colback=codebg,
    colframe=primary,
    coltitle=white,
    fonttitle=\bfseries,
    title=#1,
    arc=3mm,
    boxrule=1pt,
    left=6pt, right=6pt, top=6pt, bottom=6pt
}

\title{\textbf{\LARGE Finite-Deformation Macroscopic Stress Relaxation for Dislocation Core Simulations}\\[0.4em]
\large Kinematically Consistent Upper-Triangular Formalism, Spatial Tangent Push-Forward, and Nested Newton-Raphson Solvers}
\author{\textbf{Umut Salman} \\ CNRS / LSPM, Université Sorbonne Paris Nord, France}
\date{\today}

\begin{document}

\maketitle

\begin{abstract}
We rigorously define the cumulative macroscopic deformation gradient tensor $\bar{\mathbf{F}} \in \mathrm{GL}^+(2)$, distinguishing it from the local element-level deformation gradients $\mathbf{F}^{(e)}$. To preserve the horizontal glide line geometry without artificial crystallographic rotation, we parameterize $\bar{\mathbf{F}}$ as an upper-triangular tensor $\bar{\mathbf{F}} = \begin{pmatrix}\lambda_x & \gamma \\ 0 & \lambda_y\end{pmatrix}$ with $\bar{F}_{21} \equiv 0$, and show why the rotational gauge constraint is required to keep the Newton-Raphson stiffness nonsingular. The macroscopic residual Cauchy stress is iteratively relaxed to machine zero via a quadratic Newton-Raphson scheme driven by the exact fourth-order spatial tangent stiffness $\mathbb{C}^{\text{spatial}}(\bar{\mathbf{F}})$ computed via tensor contraction and push-forward of the internal acoustic tensor. We formulate the correct algorithmic sequence: topological Delaunay remeshing on the initial dislocated configuration, followed by normal core relaxation and full Cauchy tensor macroscopic stress relaxation, ensuring no disruptive trailing remeshing.
\end{abstract}

\vspace{0.5em}
\tableofcontents
\vspace{1em}
\hrule
\vspace{1.5em}

% ==============================================================================
\section{Introduction and Physical Problem}
% ==============================================================================

In computational continuum mechanics and atomistic-to-continuum multiscale modeling of crystalline defects, an edge dislocation is conventionally introduced into a finite domain $\Omega \subset \mathbb{R}^2$ by prescribing the analytical linear elastic displacement field $\mathbf{u}_0(\mathbf{X})$ (e.g., isotropic Volterra or anisotropic Stroh/Eshelby solution) across all atomic positions $\mathbf{X}_i \in \Omega$.

To isolate the nonlinear dislocation core from the outer boundary, the simulation cylinder is partitioned into two regions:
\begin{enumerate}
    \item \textbf{Internal Core Region ($\|\mathbf{X}_i - \mathbf{X}_{\text{core}}\| \le R_{\text{free}}$)}: 
    Nodes are free to displace under energy minimization to resolve the discrete lattice core structure, nonlinear stacking faults, and lattice discreteness effects.
    \item \textbf{Frozen Boundary Mantle ($\|\mathbf{X}_i - \mathbf{X}_{\text{core}}\| > R_{\text{free}}$)}: 
    Nodes are held fixed under Dirichlet boundary conditions $\mathbf{x}_i = \mathbf{X}_i + \mathbf{u}_0(\mathbf{X}_i)$.
\end{enumerate}

\begin{infobox}[The Finite-Size Boundary Reaction Dilemma]
While the analytical displacement field $\mathbf{u}_0$ satisfies traction-free conditions at infinity ($r \to \infty$), at any \emph{finite} boundary distance $r = L/2$, the local stress field $\sigma_{ij}^{\text{analytical}} \propto \frac{\mu b}{2\pi(1-\nu)r}$ is non-zero. Freezing the mantle nodes prevents the finite body from undergoing the macroscopic strain necessary to accommodate the extra atomic half-plane. 

Consequently, the crystal develops a macroscopic Cauchy stress tensor $\bar{\boldsymbol{\sigma}} \ne \mathbf{0}$, which distorts the core structure, shifts the apparent core energy, and creates spurious image forces that repel or attract the dislocation.
\end{infobox}

The objective of \textbf{Full Cauchy Tensor Macroscopic Stress Relaxation} is to determine the macroscopic deformation gradient $\bar{\mathbf{F}} \in \mathrm{GL}^+(2)$ such that the volume-averaged Cauchy stress tensor vanishes identically:
\begin{equation}
\bar{\boldsymbol{\sigma}}(\bar{\mathbf{F}}) = \frac{1}{\Omega(\bar{\mathbf{F}})} \int_{\Omega(\bar{\mathbf{F}})} \boldsymbol{\sigma}(\mathbf{x}) \, \mathrm{d}\Omega \longrightarrow \mathbf{0}.
\end{equation}

% ==============================================================================
\section{The Macroscopic Deformation Gradient \texorpdfstring{$\bar{\mathbf{F}}$}{F\_bar}: Definition and Kinematics}
% ==============================================================================

\subsection{Multiscale Definition: Average Deformation Gradient}

In multiscale continuum mechanics and homogenization theory, a clear distinction is maintained between:
\begin{enumerate}
    \item \textbf{Microscopic Deformation Gradient $\mathbf{F}^{(e)}(\mathbf{X})$}:
    The local gradient evaluated element-by-element inside each triangular simplex $e \in \mathcal{E}$:
    \begin{equation}
    \mathbf{F}^{(e)} = \nabla_{\mathbf{X}} \mathbf{x}^{(e)} = \mathbf{x}_e \cdot \nabla_{\mathbf{X}} \mathbf{N}^{(e)}.
    \end{equation}
    Near the dislocation core, $\mathbf{F}^{(e)}$ varies wildly across atomic spacings to capture the discrete lattice dislocation core and atomic bond distortions.
    \item \textbf{Macroscopic Deformation Gradient $\bar{\mathbf{F}}$}:
    A single, spatially uniform $2 \times 2$ second-order tensor that describes the overall macroscopic stretch, dilation, and shear of the entire crystal cylinder relative to the initial dislocated configuration. By the average deformation theorem (Stokes' theorem in the reference configuration $\Omega_0$):
    \begin{equation}
    \label{eq:F_bar_def}
    \bar{\mathbf{F}} \equiv \langle \mathbf{F} \rangle_{\Omega_0} = \frac{1}{|\Omega_0|} \int_{\Omega_0} \mathbf{F}(\mathbf{X}) \, \mathrm{d}\Omega_0 = \frac{1}{|\Omega_0|} \oint_{\partial \Omega_0} \mathbf{x} \otimes \mathbf{N}_0 \, \mathrm{d}S_0,
    \end{equation}
    where $\mathbf{N}_0$ is the outward reference unit normal along the domain boundary $\partial \Omega_0$.
\end{enumerate}

When an affine boundary condition $\mathbf{x}(\mathbf{X}) = \bar{\mathbf{F}} \mathbf{X}$ is applied to the boundary of a domain, the volume average $\langle \mathbf{F} \rangle$ is identically equal to $\bar{\mathbf{F}}$. In our code, $\bar{\mathbf{F}}$ is the primary macroscopic variable: it corresponds to \texttt{Eigen::Matrix2d F\_bar} in \texttt{src/experiments/dislocation\_study.cpp}, is exported as \texttt{"final\_F"} in \texttt{dislocation\_profiles.json}, and is surfaced in the GUI as the \textbf{Macro Stretch $\mathbf{F}$} KPI.

\subsection{Multiplicative Accumulation Across Newton Iterations}

At the start of the simulation, the analytical Volterra or Stroh displacement field $\mathbf{u}_0(\mathbf{X})$ is introduced, placing atoms at:
\begin{equation}
\mathbf{x}_i^{(0)} = \mathbf{X}_i + \mathbf{u}_0(\mathbf{X}_i - \mathbf{x}_{\text{core}}).
\end{equation}
At this initial state ($k = 0$), no macroscopic boundary adjustment has yet occurred, which corresponds to:
\begin{equation}
\bar{\mathbf{F}}^{(0)} = \mathbf{I} = \begin{pmatrix} 1 & 0 \\ 0 & 1 \end{pmatrix}.
\end{equation}
At each outer macroscopic iteration $k$, the Newton solver computes an incremental strain $\Delta\boldsymbol{\varepsilon}^{(k)}$, which constructs an incremental deformation gradient $\mathbf{F}_{\text{inc}}^{(k)}$. The cumulative macroscopic deformation gradient $\bar{\mathbf{F}}$ updates multiplicatively:
\begin{equation}
\label{eq:F_bar_update}
\bar{\mathbf{F}}^{(k+1)} = \mathbf{F}_{\text{inc}}^{(k)} \, \bar{\mathbf{F}}^{(k)} = \mathbf{F}_{\text{inc}}^{(k)} \, \mathbf{F}_{\text{inc}}^{(k-1)} \cdots \mathbf{F}_{\text{inc}}^{(0)} \, \mathbf{I}.
\end{equation}
All atomic coordinates transform relative to the dislocation core $\mathbf{x}_{\text{core}}$. For the frozen boundary mantle nodes ($\mathbf{x}_i \in \Omega_{\text{mantle}}$, $\|\mathbf{X}_i - \mathbf{x}_{\text{core}}\| > R_{\text{free}}$), their position at step $k+1$ is given by the cumulative transformation:
\begin{equation}
\label{eq:coord_cumul}
\mathbf{x}_i^{(k+1)} = \mathbf{x}_{\text{core}} + \bar{\mathbf{F}}^{(k+1)} \left( \mathbf{x}_i^{(0)} - \mathbf{x}_{\text{core}} \right) \qquad (\forall \mathbf{x}_i \in \Omega_{\text{mantle}}).
\end{equation}
For the interior free nodes ($\mathbf{x}_j \in \Omega_{\text{core}}$, $\|\mathbf{X}_j - \mathbf{x}_{\text{core}}\| \le R_{\text{free}}$), which have relaxed to an equilibrium configuration $\mathbf{x}_j^{(k)}$ under L-BFGS at step $k$, the incremental affine transformation:
\begin{equation}
\label{eq:coord_inc}
\mathbf{x}_j^{(k+1)} = \mathbf{x}_{\text{core}} + \mathbf{F}_{\text{inc}}^{(k)} \left( \mathbf{x}_j^{(k)} - \mathbf{x}_{\text{core}} \right)
\end{equation}
carries forward the relaxed core geometry as a warm start for the inner solver at iteration $k+1$. Because matrix multiplication is associative and $\mathbf{x}_i^{(k)} = \mathbf{x}_{\text{core}} + \bar{\mathbf{F}}^{(k)} (\mathbf{x}_i^{(0)} - \mathbf{x}_{\text{core}})$ for the unrelaxed mantle nodes, applying $\mathbf{F}_{\text{inc}}^{(k)}$ incrementally to all nodes (as done in the code and Eq.~\eqref{eq:coord_update}) is identically equivalent to Eq.~\eqref{eq:coord_cumul} for the mantle while preserving core relaxation.

\subsection{Component Breakdown and Physical Meaning}

In 2D Cartesian coordinates, $\bar{\mathbf{F}}$ is parameterized as:
\begin{equation}
\label{eq:F_bar_matrix}
\bar{\mathbf{F}} = \begin{pmatrix} \bar{F}_{11} & \bar{F}_{12} \\ \bar{F}_{21} & \bar{F}_{22} \end{pmatrix} = \begin{pmatrix} \lambda_x & \gamma \\ 0 & \lambda_y \end{pmatrix}.
\end{equation}
Each entry of $\bar{\mathbf{F}}$ has a precise physical and crystallographic role:
\begin{itemize}
    \item $\bar{F}_{11} = \lambda_x$: 
    \textbf{Longitudinal Stretch along the Slip Direction}. Measures the uniform expansion or compression of the crystal parallel to the Burgers vector $\mathbf{b} \parallel \hat{\mathbf{e}}_1$. Adjusting $\lambda_x$ relieves the macroscopic normal stress $\bar{\sigma}_{xx}$.
    \item $\bar{F}_{22} = \lambda_y$: 
    \textbf{Transverse Normal Stretch}. Measures the crystal dilation perpendicular to the slip plane. Because introducing an edge dislocation inserts an extra half-plane of atoms, the finite crystal must expand vertically (and adjust via Poisson contraction) to accommodate the extra material. Adjusting $\lambda_y$ relieves the macroscopic normal stress $\bar{\sigma}_{yy}$.
    \item $\bar{F}_{12} = \gamma$: 
    \textbf{Simple Shear Strain Parallel to the Glide Plane}. Relieves the macroscopic resolved shear stress $\bar{\sigma}_{xy}$. In crystalline plasticity, dislocation motion occurs via simple shear along the slip line.
    \item $\bar{F}_{21} \equiv 0$: 
    \textbf{Kinematic Constraint (Zero Crystallographic Tilt)}. Enforces that lines parallel to the $x$-axis (the dislocation glide line) do not undergo artificial rigid rotation relative to the Cartesian frame.
\end{itemize}

\subsection{Why \texorpdfstring{$\bar{F}_{21} \equiv 0$}{F\_21 = 0}: Slip Plane Invariance}

Consider an arbitrary affine deformation centered at the core $\mathbf{x}_{\text{core}} = (x_c, y_c)^T$. The deformed vertical coordinate $y'$ of any point $(x, y)$ is:
\begin{equation}
y' - y_c = \bar{F}_{21}(x - x_c) + \bar{F}_{22}(y - y_c).
\end{equation}
For any atom lying on the dislocation slip plane ($y = y_c$):
\begin{equation}
y' - y_c = \bar{F}_{21}(x - x_c).
\end{equation}
\begin{itemize}
    \item If $\bar{F}_{21} \ne 0$ (such as in a naive symmetric stretch where $\bar{F}_{21} = \bar{F}_{12} = \varepsilon_{12}$), then $y' - y_c = \varepsilon_{12}(x - x_c) \ne 0$. The entire slip plane is tilted at an angle $\theta \approx \bar{F}_{21} / \bar{F}_{22}$. This violates the physical boundary conditions: the cut line and the Burgers vector are no longer aligned with the grid!
    \item If $\bar{F}_{21} \equiv 0$ (upper-triangular parameterization), then:
    \begin{equation}
    y' - y_c = 0 \cdot (x - x_c) + \lambda_y \cdot 0 \equiv 0 \qquad \forall x \in \mathbb{R}.
    \end{equation}
    The glide plane remains \textbf{strictly horizontal and invariant} at all times.
\end{itemize}

\subsection{Why Cannot \texorpdfstring{$\bar{F}_{12}$}{F\_12} and \texorpdfstring{$\bar{F}_{21}$}{F\_21} Both Be Independent?}

A natural question arises: \emph{Why can we not let all 4 components of $\bar{\mathbf{F}}$ vary independently during stress relaxation?}

The answer stems from the fundamental balance of angular momentum and material frame indifference:
\begin{enumerate}
    \item \textbf{3 Equations vs 4 Unknowns}:
    The Cauchy stress tensor $\bar{\boldsymbol{\sigma}}$ is strictly symmetric:
    \begin{equation}
    \bar{\boldsymbol{\sigma}} = \begin{pmatrix} \bar{\sigma}_{xx} & \bar{\sigma}_{xy} \\ \bar{\sigma}_{xy} & \bar{\sigma}_{yy} \end{pmatrix} \implies \bar{\sigma}_{yx} \equiv \bar{\sigma}_{xy}.
    \end{equation}
    Therefore, setting $\bar{\boldsymbol{\sigma}} = \mathbf{0}$ yields only \textbf{3 independent scalar equations} ($\bar{\sigma}_{xx} = 0, \bar{\sigma}_{yy} = 0, \bar{\sigma}_{xy} = 0$). Solving 3 equations with 4 unconstrained variables is an underdetermined problem with an infinite family of solutions.
    \item \textbf{The 4th Degree of Freedom is Pure Rigid Rotation}:
    By polar decomposition, $\bar{\mathbf{F}} = \mathbf{R}(\theta)\mathbf{U}$. The right Cauchy-Green strain tensor $\mathbf{C} = \bar{\mathbf{F}}^T \bar{\mathbf{F}} = \mathbf{U}^T \mathbf{R}^T \mathbf{R} \mathbf{U} = \mathbf{U}^2$ completely filters out rigid-body rotation $\theta$. Because the strain energy $\Phi(\mathbf{C})$ and Cauchy stress are frame-indifferent (objective), a pure rigid rotation in space generates \textbf{zero stress change}:
    \begin{equation}
    \frac{\partial \bar{\boldsymbol{\sigma}}}{\partial \theta} \equiv \mathbf{0}.
    \end{equation}
    \item \textbf{Singular Tangent Matrix}:
    If a $4 \times 4$ Newton-Raphson tangent matrix $\frac{\partial \bar{\boldsymbol{\sigma}}}{\partial \bar{\mathbf{F}}}$ were constructed, its 4th column would be identically in the nullspace:
    \begin{equation}
    \det\left( \frac{\partial \bar{\boldsymbol{\sigma}}}{\partial \bar{\mathbf{F}}} \right) = 0.
    \end{equation}
    Newton's method would fail with a division by zero or uncontrolled rotational drift.
\end{enumerate}

Fixing $\bar{F}_{21} \equiv 0$ is the unique gauge choice that eliminates this singular rotational nullspace while simultaneously preserving the horizontal orientation of the dislocation slip plane.

% ==============================================================================
\section{Macroscopic Cauchy Stress Tensor Formulation}
% ==============================================================================

In our nonlinear finite element framework, the crystal energy is evaluated element-by-element over triangular elements $e \in \mathcal{E}$.

\subsection{Element Deformation and Metric Tensor}

For each triangular element $e$ with shape function derivatives $\nabla_{\mathbf{X}} \mathbf{N}^{(e)}$ in the reference crystal, the total deformation gradient $\mathbf{F}^{(e)}$ and right Cauchy-Green tensor $\mathbf{C}^{(e)}$ are:
\begin{align}
\mathbf{F}^{(e)} &= \mathbf{x}_e \cdot \nabla_{\mathbf{X}} \mathbf{N}^{(e)}, \\
\mathbf{C}^{(e)} &= \mathbf{F}^{(e)T} \mathbf{F}^{(e)}.
\end{align}
To respect the discrete point-group symmetry $\mathcal{G} \subset \mathrm{GL}(2, \mathbb{Z})$ of the underlying crystal lattice (e.g., cubic/square or hexagonal/triangular), $\mathbf{C}^{(e)}$ is mapped into the fundamental Lagrange-reduced domain $\mathcal{D}_L$ via an integer unimodular matrix $\mathbf{m}^{(e)} \in \mathrm{GL}(2, \mathbb{Z})$:
\begin{equation}
\mathbf{C}_{\text{red}}^{(e)} = \mathbf{m}^{(e)T} \mathbf{C}^{(e)} \mathbf{m}^{(e)}.
\end{equation}

\subsection{First Piola-Kirchhoff and Cauchy Stress Tensors}

The element strain energy density $\Phi(\mathbf{C}_{\text{red}})$ yields the second Piola-Kirchhoff stress in the reduced space:
\begin{equation}
\mathbf{S}_{\text{red}}^{(e)} = 2 \frac{\partial \Phi}{\partial \mathbf{C}_{\text{red}}^{(e)}}.
\end{equation}
Pulling back via the reduction matrix $\mathbf{Z}^{(e)} = \mathbf{m}^{(e)}$ gives the first Piola-Kirchhoff stress $\mathbf{P}^{(e)}$:
\begin{equation}
\mathbf{P}^{(e)} = 2 \mathbf{F}^{(e)} \mathbf{Z}^{(e)} \frac{\partial \Phi}{\partial \mathbf{C}_{\text{red}}^{(e)}} \mathbf{Z}^{(e)T}.
\end{equation}
The true spatial Cauchy stress $\boldsymbol{\sigma}^{(e)}$ is obtained by the standard push-forward:
\begin{equation}
\boldsymbol{\sigma}^{(e)} = \frac{1}{\det \mathbf{F}^{(e)}} \mathbf{P}^{(e)} \mathbf{F}^{(e)T}.
\end{equation}

\subsection{Volume-Averaged Macroscopic Stress}

The macroscopic Cauchy stress $\bar{\boldsymbol{\sigma}}$ driving the Newton relaxation is evaluated as the area-weighted volume average over the active elements $\mathcal{E}_{\text{active}}$:
\begin{equation}
\label{eq:macro_stress}
\bar{\boldsymbol{\sigma}} = \frac{\sum_{e \in \mathcal{E}_{\text{active}}} \boldsymbol{\sigma}^{(e)} \Omega^{(e)}}{\sum_{e \in \mathcal{E}_{\text{active}}} \Omega^{(e)}}, \qquad \Omega^{(e)} = \det(\mathbf{F}^{(e)}) \Omega_0^{(e)}.
\end{equation}
Here $\mathcal{E}_{\text{active}}$ comprises all triangular elements containing at least one free node in the core relaxation cylinder ($r \le R_{\text{free}}$). Restricting the integration domain to $\mathcal{E}_{\text{active}}$ directly isolates the non-vanishing reaction stresses generated by the relaxed atomic core, preventing them from being diluted by the fixed analytical displacement field prescribed on the frozen Dirichlet mantle ($r > R_{\text{free}}$).
Due to the symmetry of $\mathbf{C}$, $\bar{\boldsymbol{\sigma}}$ is strictly symmetric: $\bar{\sigma}_{xy} = \bar{\sigma}_{yx}$. In Voigt notation:
\begin{equation}
\bar{\boldsymbol{\sigma}}_{\text{voigt}} = \begin{pmatrix} \bar{\sigma}_{xx} \\ \bar{\sigma}_{yy} \\ \bar{\sigma}_{xy} \end{pmatrix} \in \mathbb{R}^3.
\end{equation}

% ==============================================================================
\section{Instantaneous Spatial Tangent Stiffness (\texorpdfstring{$\mathbb{C}^{\text{spatial}}$}{C\_spatial})}
% ==============================================================================

To achieve true quadratic Newton convergence when solving $\bar{\boldsymbol{\sigma}}(\bar{\mathbf{F}}) = \mathbf{0}$, we require the exact derivative of Cauchy stress with respect to macroscopic strain:
\begin{equation}
\mathrm{d}\bar{\boldsymbol{\sigma}} = \mathbb{C}^{\text{spatial}} : \mathrm{d}\boldsymbol{\varepsilon}.
\end{equation}

In our code, this is evaluated directly by \texttt{AcousticTensor::getSpatialTangent()} in \texttt{src/acoustic\_tensor.cpp}.

\subsection{Material Tangent Computation}

At the current macroscopic state $\bar{\mathbf{F}}$, the acoustic tensor evaluates the fourth-order material tangent modulus $\mathbb{A}_{KLMN} = \frac{\partial^2 \Phi}{\partial F_{kM} \partial F_{lN}}$ using the second derivatives $\frac{\partial^2 \Phi}{\partial \mathbf{C}^2}$ calculated analytically by \texttt{BaseLatticeCalculator}.

\subsection{Spatial Push-Forward and Geometric Stress Term}

The push-forward to the spatial configuration consists of two essential contributions:
\begin{equation}
\mathbb{C}_{ijkl}^{\text{spatial}} = \underbrace{\frac{1}{J} \bar{F}_{iK} \bar{F}_{jL} \bar{F}_{kM} \bar{F}_{lN} \mathbb{A}_{KLMN}}_{\text{Term 1: Material Push-Forward}} + \underbrace{\frac{1}{2}\left( \bar{\sigma}_{ik}\delta_{jl} + \bar{\sigma}_{jk}\delta_{il} + \bar{\sigma}_{il}\delta_{jk} + \bar{\sigma}_{jl}\delta_{ik} \right)}_{\text{Term 2: Geometric Cauchy Stress Term}},
\end{equation}
where $J = \det(\bar{\mathbf{F}})$. The geometric stress term ensures rigorous consistency with the Jaumann rate of Cauchy stress under finite deformation.

\begin{infobox}[Spatial Tangent Behavior Near Equilibrium]
While the exact directional derivative of Cauchy stress $\bar{\boldsymbol{\sigma}} = J^{-1} \mathbf{P}\bar{\mathbf{F}}^T$ with respect to an upper-triangular increment $\Delta \mathbf{L} = \mathbf{F}_{\text{inc}} - \mathbf{I}$ differs slightly from the symmetric Jaumann tangent by geometric terms proportional to $\bar{\boldsymbol{\sigma}}$ (e.g., $-\bar{\boldsymbol{\sigma}}\,\mathrm{tr}(\Delta\mathbf{L})$ and convection terms), all terms linear in $\bar{\boldsymbol{\sigma}}$ vanish identically as $\bar{\boldsymbol{\sigma}} \to \mathbf{0}$. Consequently, near the stress-relaxed state, $\mathbb{C}^{\text{spatial}}$ asymptotically coincides with the exact Jacobian of the macroscopic residual stress vector, guaranteeing robust quadratic Newton convergence.
\end{infobox}

\subsection{Voigt Matrix Representation}

The fourth-order spatial tensor $\mathbb{C}_{ijkl}^{\text{spatial}}$ is mapped into a $3 \times 3$ matrix $\mathbf{C}_{\text{voigt}}$:
\begin{equation}
\label{eq:C_voigt}
\mathbf{C}_{\text{voigt}} = \begin{pmatrix}
\mathbb{C}_{1111} & \mathbb{C}_{1122} & \mathbb{C}_{1112} \\
\mathbb{C}_{2211} & \mathbb{C}_{2222} & \mathbb{C}_{2212} \\
\mathbb{C}_{1211} & \mathbb{C}_{1222} & \mathbb{C}_{1212}
\end{pmatrix}.
\end{equation}

% ==============================================================================
\section{The Newton-Raphson Inversion and Update Algorithm}
% ==============================================================================

At macroscopic iteration $k$, given the current residual stress $\bar{\boldsymbol{\sigma}}^{(k)}$ and tangent stiffness $\mathbf{C}_{\text{voigt}}^{(k)}$, we seek the strain increment:
\begin{equation}
\Delta\boldsymbol{\varepsilon} = \begin{pmatrix} \Delta\varepsilon_{11} \\ \Delta\varepsilon_{22} \\ 2\Delta\varepsilon_{12} \end{pmatrix}.
\end{equation}

\subsection{Solving the Linearized System}

We solve:
\begin{equation}
\mathbf{C}_{\text{voigt}}^{(k)} \Delta\boldsymbol{\varepsilon} = -\bar{\boldsymbol{\sigma}}^{(k)}.
\end{equation}
Using Eigen's full-pivoting LU decomposition (\texttt{fullPivLu}):
\begin{equation}
\Delta\boldsymbol{\varepsilon} = - \left( \mathbf{C}_{\text{voigt}}^{(k)} \right)^{-1} \bar{\boldsymbol{\sigma}}^{(k)}.
\end{equation}
If $\det(\mathbf{C}_{\text{voigt}}) < 10^{-10}$ (near elastic instability or phase transition), a singular value decomposition (SVD) pseudoinverse is automatically invoked as a robust fallback.

\subsection{Step Size Safeguard}

To ensure unconditional stability even far from equilibrium, the strain increment is clamped:
\begin{equation}
\Delta\varepsilon_i \longleftarrow \mathrm{sign}(\Delta\varepsilon_i) \min\left( |\Delta\varepsilon_i|, \, \Delta\varepsilon_{\max} \right), \qquad \Delta\varepsilon_{\max} = 0.02.
\end{equation}

\subsection{Upper-Triangular Incremental Update}

The incremental deformation gradient $\mathbf{F}_{\text{inc}}$ is constructed directly in upper-triangular form:
\begin{equation}
\label{eq:F_inc}
\mathbf{F}_{\text{inc}} = \begin{pmatrix}
1 + \Delta\varepsilon_{11} & 2\Delta\varepsilon_{12} \\
0 & 1 + \Delta\varepsilon_{22}
\end{pmatrix}.
\end{equation}
The cumulative macroscopic deformation gradient updates multiplicatively:
\begin{equation}
\bar{\mathbf{F}}^{(k+1)} = \mathbf{F}_{\text{inc}} \, \bar{\mathbf{F}}^{(k)}.
\end{equation}
Since both $\mathbf{F}_{\text{inc}}$ and $\bar{\mathbf{F}}^{(k)}$ have vanishing $(2,1)$ components, their product satisfies:
\begin{equation}
(\bar{\mathbf{F}}^{(k+1)})_{21} = 0 \cdot \bar{F}_{11} + (1 + \Delta\varepsilon_{22}) \cdot 0 \equiv 0.
\end{equation}
The upper-triangular structure is preserved \emph{to machine precision at all steps}.

\subsection{Core-Centered Atomic Coordinate Shift}

All atomic coordinates $\mathbf{x}_j$ (both interior free nodes and boundary mantle nodes) are updated by an affine expansion \emph{centered at the dislocation core} $\mathbf{x}_{\text{core}} = (x_c, y_c)^T$:
\begin{equation}
\label{eq:coord_update}
\mathbf{x}_j^{(k+1)} = \mathbf{x}_{\text{core}} + \mathbf{F}_{\text{inc}} \left( \mathbf{x}_j^{(k)} - \mathbf{x}_{\text{core}} \right).
\end{equation}

\begin{infobox}[Why Centering at $\mathbf{x}_{\text{core}}$ is Critical]
If coordinates were scaled relative to the domain origin $(0, 0)$, any non-zero shear $\Delta\gamma = 2\Delta\varepsilon_{12}$ or normal strain would translate the core position:
$$\mathbf{x}_{\text{core}} \longrightarrow \mathbf{F}_{\text{inc}} \mathbf{x}_{\text{core}} \ne \mathbf{x}_{\text{core}}.$$
This translation would misalign the core with the cylindrical relaxation window $R_{\text{free}}$ and introduce artificial core drift. Centering the transformation at $\mathbf{x}_{\text{core}}$ guarantees that $\mathbf{x}_{\text{core}} - \mathbf{x}_{\text{core}} = \mathbf{0}$, anchoring the core invariant in space.
\end{infobox}

% ==============================================================================
\section{Algorithmic Ordering and Nested Optimization Architecture}
% ==============================================================================

The dislocation relaxation algorithm operates in three strictly ordered phases to guarantee both topological compatibility and macroscopic zero-stress equilibrium:
\begin{enumerate}
    \item \textbf{Phase 0: Initial Mesh Generation (Delaunay Remeshing on Dislocated State)}:
    If \texttt{enable\_remeshing} is active, the 2D Delaunay triangulation is constructed directly from the dislocated atom coordinates $\mathbf{x}_i^{(0)} = \mathbf{X}_i + \mathbf{u}_0(\mathbf{X}_i - \mathbf{x}_{\text{core}})$. This dynamically reconnects bonds across the displacement cut and accommodates dislocation core dilation. If remeshing is disabled, the fixed pristine lattice connectivity is used.
    \item \textbf{Phase 1: Normal Core Relaxation (Fixed Dirichlet Mantle)}:
    Holding the outer mantle nodes ($\|\mathbf{X}_i - \mathbf{x}_{\text{core}}\| > R_{\text{free}}$) fixed under the initial boundary conditions ($\bar{\mathbf{F}} = \mathbf{I}$), L-BFGS minimizes the energy with respect to the coordinates of the free core nodes ($\|\mathbf{X}_i - \mathbf{x}_{\text{core}}\| \le R_{\text{free}}$).
    \item \textbf{Phase 2: Full Cauchy Tensor Macro-Stress Relaxation}:
    If \texttt{relax\_macro\_stress} is enabled, an outer Newton loop iteratively updates the macroscopic deformation gradient $\bar{\mathbf{F}}$ using the spatial tangent stiffness $\mathbb{C}^{\text{spatial}}(\bar{\mathbf{F}})$ from the acoustic tensor. At each macro step, all coordinates are scaled relative to $\mathbf{x}_{\text{core}}$ and internal free nodes are re-relaxed until $\max|\bar{\sigma}_{ij}| < \text{stress\_tol}$.
    \item \textbf{No Trailing Remeshing}:
    Once the macro-stress relaxation converges, \emph{no further remeshing is performed}. Retriangulating after Cauchy relaxation would alter local element references and reintroduce spurious residual stresses, destroying the zero-stress macroscopic equilibrium.
\end{enumerate}

\begin{algorithm}[h!]
\caption{Full Dislocation Relaxation: Remesh $\to$ Normal Relaxation $\to$ Cauchy Relaxation}
\label{alg:macro_stress}
\begin{algorithmic}[1]
\Require Reference Lattice $\{ \mathbf{X}_i \}_{i=1}^N$, Core $\mathbf{x}_{\text{core}}$, Radius $R_{\text{free}}$, Tolerance $\eta$, Max Macro Iters $K_{\max}$, \texttt{enable\_remeshing}, \texttt{relax\_macro\_stress}
\State Apply analytical Volterra / Stroh displacement: $\mathbf{x}_i^{(0)} \gets \mathbf{X}_i + \mathbf{u}_0(\mathbf{X}_i - \mathbf{x}_{\text{core}})$
\State Partition nodes into free ($\|\mathbf{X}_i - \mathbf{x}_{\text{core}}\| \le R_{\text{free}}$) and Dirichlet boundary ($> R_{\text{free}}$)
\If{\texttt{enable\_remeshing}}
    \State Build Delaunay triangulation from dislocated positions $\mathbf{x}^{(0)}$ (with core-edge filtering)
\Else
    \State Build triangulation from reference lattice positions $\mathbf{X}$
\EndIf
\State \textbf{Phase 1 (Normal Relaxation)}: Relax free core nodes under fixed mantle: $\mathbf{x}_{\text{free}} \gets \arg\min E(\mathbf{x}_{\text{free}} \,|\, \mathbf{x}_{\text{bound}})$
\If{\texttt{relax\_macro\_stress}}
    \State \textbf{Phase 2 (Macro-Stress Relaxation)}: Initialize $\bar{\mathbf{F}} \gets \mathbf{I}$, converged $\gets \mathbf{false}$
    \For{$k = 0$ \textbf{to} $K_{\max} - 1$}
        \State Compute volume-averaged Cauchy stress $\bar{\boldsymbol{\sigma}}$ via Eq.~\eqref{eq:macro_stress}
        \State $\sigma_{\max} \gets \max\left( |\bar{\sigma}_{xx}|, \, |\bar{\sigma}_{yy}|, \, |\bar{\sigma}_{xy}| \right)$
        \If{$\sigma_{\max} < \eta$}
            \State converged $\gets \mathbf{true}$; \textbf{break}
        \EndIf
        \State Compute $\mathbb{C}^{\text{spatial}}(\bar{\mathbf{F}})$ via \texttt{AcousticTensor::getSpatialTangent()}
        \State Extract $3 \times 3$ Voigt matrix $\mathbf{C}_{\text{voigt}}$ via Eq.~\eqref{eq:C_voigt}
        \State Solve $\mathbf{C}_{\text{voigt}} \Delta\boldsymbol{\varepsilon} = -\bar{\boldsymbol{\sigma}}_{\text{voigt}}$
        \State Clamp $\Delta\varepsilon_i \gets \mathrm{clamp}(\Delta\varepsilon_i, -0.02, 0.02)$
        \State Construct upper-triangular $\mathbf{F}_{\text{inc}}$ via Eq.~\eqref{eq:F_inc}
        \State Update cumulative deformation: $\bar{\mathbf{F}} \gets \mathbf{F}_{\text{inc}} \bar{\mathbf{F}}$
        \State Update all atom positions: $\mathbf{x}_j \gets \mathbf{x}_{\text{core}} + \mathbf{F}_{\text{inc}} (\mathbf{x}_j - \mathbf{x}_{\text{core}})$
        \State Synchronize solver array and update element reference mesh
        \State Re-relax internal free nodes: $\mathbf{x}_{\text{free}} \gets \arg\min E(\mathbf{x}_{\text{free}} \,|\, \mathbf{x}_{\text{bound}})$
    \EndFor
\EndIf
\State Export converged relaxed state to VTK, CSV, and JSON (without trailing remeshing)
\end{algorithmic}
\end{algorithm}

% ==============================================================================
\section{Numerical Verification and Finite-Size Scaling}
% ==============================================================================

\subsection{Quadratic Convergence Demonstration}

We illustrate the convergence behavior on a $25 \times 25$ square lattice with a central edge dislocation ($a = 1.0$, $R_{\text{free}} = 8.0$, Zanzotto continuum potential $\beta = -0.25, K = 5.0$):
\begin{itemize}
    \item \textbf{Step 0 (Initial unrelaxed mantle, $\bar{\mathbf{F}}^{(0)} = \mathbf{I}$)}:
    \begin{align*}
    E &= 1.2041720450 \\
    \bar{\boldsymbol{\sigma}} &= \begin{pmatrix} -0.051077 & -0.000020 \\ -0.000020 & -0.045825 \end{pmatrix}, \quad \max|\bar{\sigma}| = 5.11 \times 10^{-2} \\
    \Delta\boldsymbol{\varepsilon} &= (1.97 \times 10^{-3}, \, 4.51 \times 10^{-4}, \, 3.23 \times 10^{-6})^T
    \end{align*}
    \item \textbf{Step 1 (After 1 Newton step, updated $\bar{\mathbf{F}}^{(1)}$)}:
    \begin{align*}
    E &= 1.1904803911 \quad (\Delta E = -0.01369) \\
    \bar{\boldsymbol{\sigma}} &= \begin{pmatrix} -4.68 \times 10^{-5} & -8.63 \times 10^{-6} \\ -8.63 \times 10^{-6} & -3.38 \times 10^{-5} \end{pmatrix}, \quad \max|\bar{\sigma}| = 4.69 \times 10^{-5}
    \end{align*}
    \textbf{Convergence reached in 2 iterations!} The maximum stress dropped by over 1,000$\times$ in a single step, confirming second-order quadratic Newton convergence.
\end{itemize}

\subsection{Finite-Size Scaling: Why \texorpdfstring{$100 \times 100$}{100x100} Behaves Differently}

A key physical insight arises when comparing small ($25 \times 25$) versus large ($100 \times 100$) systems:
\begin{table}[h!]
\centering
\small
\setlength{\tabcolsep}{6pt}
\begin{tabular}{lcccc}
\toprule
\textbf{Grid Size} & \textbf{Atoms} & \textbf{Unrelaxed $\max|\bar{\sigma}|$} & \textbf{Relaxed $\max|\bar{\sigma}|$} & \textbf{Box Expansion $\Delta L / L$} \\
\midrule
$25 \times 25$   & 625    & $5.11 \times 10^{-2}$  & $4.69 \times 10^{-5}$  & $+0.20\%$ \\
$100 \times 100$ & 10,000 & $4.40 \times 10^{-4}$  & $9.38 \times 10^{-14}$ & $+0.007\%$ \\
\bottomrule
\end{tabular}
\caption{Comparison of boundary reaction stress and macroscopic relaxation across system sizes.}
\end{table}

\begin{infobox}[Nonlinear Anharmonicity and Core Dilatation Scaling]
A striking observation from Table 1 is that when $L$ increases by $4\times$ (from $25 \times 25$ to $100 \times 100$), the unrelaxed reaction stress $\max|\bar{\boldsymbol{\sigma}}|$ drops by a factor of \textbf{116}, decaying much faster than the nominal linear-elastic traction rate $\propto 1/L$.

This rapid decay has a profound physical origin:
\begin{enumerate}
    \item \textbf{Antisymmetry of Linear Elastic Stress}: In classical linear elasticity, the Volterra/Stroh stress tensor $\boldsymbol{\sigma}^{\text{linear}} \propto \frac{\mu b}{r}$ is antisymmetric across the slip plane ($y \leftrightarrow -y$). When integrated over a symmetric crystal domain centered at the core, the linear elastic contribution vanishes identically:
    $$\int_{\Omega} \boldsymbol{\sigma}^{\text{linear}}(\mathbf{x}) \, \mathrm{d}\Omega \equiv \mathbf{0}.$$
    \item \textbf{Nonlinear Elastic Anharmonicity}: Consequently, the non-vanishing net macroscopic stress $\bar{\boldsymbol{\sigma}}$ is entirely driven by higher-order nonlinear elasticity ($\sigma \propto \varepsilon^2 \sim 1/r^2$) and the net insertion volume of the core defect ($\Delta V \sim b^2$).
    \item \textbf{Geometric Volume Dilution}: The defect insertion strain scales with the domain area as $\Delta V / |\Omega| \propto b^2 / L^2$. Together with nonlinear core accommodation, this produces an effective scaling $\bar{\sigma} \sim \mathcal{O}(L^{-2})$ to $\mathcal{O}(L^{-3})$.
\end{enumerate}
This confirms that the macroscopic reaction stress in a dislocated crystal is intrinsically a \emph{nonlinear finite-deformation phenomenon}, providing a rigorous physical justification for our finite-deformation formulation over standard linear superposition.
\end{infobox}

In a $100 \times 100$ crystal ($L = 100$), the initial residual stress is already very small ($4.4 \times 10^{-4}$), and the macroscopic strain required to relieve it is merely $\sim 0.007\%$, shifting outer boundaries by less than $0.007$ lattice units. Consequently, the local core fitting parameters ($E_{\text{core}}, K$) are virtually unaffected by the distant boundary at $L = 100$. For smaller or moderate domain sizes ($L \le 30$), however, macroscopic relaxation is vital to prevent severe artificial core compression.

% ==============================================================================
\section{Code Mapping and User Guide}
% ==============================================================================

The implementation spans three tightly integrated layers:
\begin{enumerate}
    \item \textbf{C++ Physics Engine}:
    \begin{itemize}
        \item \texttt{include/experiments/dislocation\_study.h}: Added \texttt{relax\_macro\_stress}, \texttt{stress\_tol}, and \texttt{max\_macro\_iters}.
        \item \texttt{src/experiments/dislocation\_study.cpp}: Implemented the macro-stress Newton loop, spatial tangent extraction, upper-triangular update $\bar{\mathbf{F}} \gets \mathbf{F}_{\text{inc}}\bar{\mathbf{F}}$, and core-centered coordinate scaling.
        \item \texttt{src/main.cpp}: Added CLI flags \texttt{--relax-stress=1}, \texttt{--stress-tol=1e-5}, and \texttt{--max-macro-iters=30}.
    \end{itemize}
    \item \textbf{Data Logging}:
    \begin{itemize}
        \item \texttt{dislocation\_profiles.json}: Exports \texttt{macro\_stress\_relaxed}, \texttt{final\_F} ($\equiv \bar{\mathbf{F}}$), and \texttt{final\_macro\_stress}.
    \end{itemize}
    \item \textbf{Lattice Launcher macOS GUI}:
    \begin{itemize}
        \item \texttt{SimulationConfig.swift}: Added toggle, fields, and automatic \texttt{\_relax} folder tagging.
        \item \texttt{ContentView.swift}: Interactive UI toggle under \emph{Mode \& Triangulation}.
        \item \texttt{DislocationProfileChartView.swift}: Dedicated KPI pills for \textbf{Macro Cauchy Stress} (showing $\max|\sigma|$ with green/red status) and \textbf{Macro Stretch $\mathbf{F}$} (showing $\bar{\mathbf{F}}$).
    \end{itemize}
\end{enumerate}

\end{document}
